\documentclass[11pt]{article}
\usepackage[margin=1in]{geometry}
\usepackage[T1]{fontenc}
\usepackage{lmodern}
\emergencystretch=2em
\usepackage{amsmath,amssymb,amsthm}
\usepackage[hidelinks]{hyperref}
\newtheorem{theorem}{Theorem}
\newtheorem{lemma}{Lemma}
\newcommand{\E}{\mathbb E}
\newcommand{\R}{\mathbb R}
\newcommand{\Pp}{\mathbb P}
\title{Disproof of Conjecture 00000007788\\[3pt]\large A lower bound for the missing area of a random polygon}
\author{C0ldSmi1e}
\date{}
\begin{document}
\maketitle
\begin{abstract}
The conjecture predicts an expected volume deficit of order
$N^{-2/(n-1)}$ for the convex hull of $N$ independent points uniform in
the Euclidean ball. In dimension two, we prove that the deficit is at
least $1/[4(N+1)]$. Its quadratic normalization therefore tends to
$+\infty$, contradicting the asserted inverse-square asymptotic for
every finite leading constant. The accompanying Lean project formalizes
the probability law, actual convex-hull area, measurability,
integrability, quantitative bound, and asymptotic contradiction.
\end{abstract}

\section{Statement and scope}
Both language versions define $K_N=\operatorname{conv}(X_1,\ldots,X_N)$,
where the $X_i$ are independent and uniform on $K$. For a solid
Euclidean ball, this is normalized volume on the ball. The later
reference to spherical measure does not change that sampling law.
Dimension $n=2$ is an admissible instance of the displayed exponent.

Let $B=\{x\in\R^2:\lVert x\rVert\leq1\}$, let $\lambda$ be Lebesgue
area, and set
\[
 \mu(A)=\frac{\lambda(A\cap B)}{\pi},\qquad
 H_N(x)=\operatorname{conv}\{x_i:1\leq i\leq N\},\qquad
 D_N=1-\frac{\E_{\mu^{\otimes N}}\lambda(H_N)}{\pi}.
\]
The necessary ball assertion of the conjecture gives
$D_N\sim c_2N^{-2}$ for a finite positive constant $c_2$.
Refuting this necessary assertion disproves the original conjunction;
it requires no choice of definitions for its additional constant-ratio
or simplex corner-measure clauses.

\begin{theorem}\label{thm:main}
For every integer $N\geq1$,
\[
 D_N\geq\frac{1}{4(N+1)}.
\]
Consequently $N^2D_N\to+\infty$. In particular, $N^2D_N$ has no finite
real limit, and $D_N\not\sim cN^{-2}$ for every $c>0$.
\end{theorem}

\newpage
\section{Measurability and existence of the expectation}
For completeness, the expectation above is an ordinary expectation of
an integrable random variable. For each fixed $N$, put
\[
 \Delta_N=\{w\in\R^N:w_i\geq0,\ \textstyle\sum_iw_i=1\}.
\]
This simplex is compact, and
$H_N(x)=\{\sum_iw_ix_i:w\in\Delta_N\}$.
The set of $(x,y,w)$ with $w\in\Delta_N$ and
$y=\sum_iw_ix_i$ is closed by continuity of the finite sum.
Projection along the compact factor $\Delta_N$ is a closed map, so
\[
 G_N=\{(x,y):y\in H_N(x)\}
\]
is closed. The measurable-section theorem for product measurable
spaces shows that $x\mapsto\lambda(H_N(x))$ is measurable.
Each finite hull is compact and has finite area. Under
$\mu^{\otimes N}$ all sample points lie in $B$ almost surely, and
convexity gives $H_N\subseteq B$. Thus
$0\leq\lambda(H_N)/\pi\leq1$ almost surely. Both normalized and
unnormalized hull area are integrable. The formalization proves these
facts explicitly before using integral monotonicity.

\section{The elementary lower bound}
For $0\leq r\leq1$, let
$E_r=\{x:\lVert x_i\rVert\leq r\text{ for all }i\}$.
Area scaling and independence give
\[
 \mu(rB)=r^2,\qquad \Pp(E_r)=(r^2)^N.
\]
On $E_r$, convexity gives $H_N\subseteq rB$, hence
$\lambda(H_N)/\pi\leq r^2$. Off this event the almost-sure bound by
$1$ remains valid. Therefore
\[
 \frac{\lambda(H_N)}{\pi}
 \leq1-(1-r^2)1_{E_r}
 \quad\text{almost surely}.
\]
Taking expectations and rearranging proves
\begin{equation}\label{eq:radius}
 D_N\geq(1-r^2)(r^2)^N.
\end{equation}
Choose
\[
 t=\frac{1}{2(N+1)},\qquad r=\sqrt{1-t}.
\]
Here $0<t\leq1/2$ and $0\leq r\leq1$. Bernoulli's inequality yields
\[
 (r^2)^N=(1-t)^N\geq1-Nt\geq\frac12.
\]
Substituting into \eqref{eq:radius} gives
$D_N\geq t/2=1/[4(N+1)]$.
For $N\geq1$ it follows that
\[
 N^2D_N\geq\frac{N^2}{4(N+1)}\geq\frac N8\longrightarrow+\infty.
\]
For each $c>0$, the ratio $D_N/(cN^{-2})=N^2D_N/c$ also tends to
$+\infty$, and cannot tend to $1$. This proves
Theorem~\ref{thm:main}. The bound is sufficient for a disproof and makes
no claim about the sharp asymptotic order of the deficit.

\newpage
\section{Correspondence with the Lean project}
The project uses Lean 4.19.0 and Mathlib v4.19.0, pinned by the supplied
toolchain and dependency manifest. The mathematical modules are
\nolinkurl{Conjecture7788.lean},
\nolinkurl{Conjecture7788/HullMeasurable.lean}, and
\nolinkurl{Conjecture7788/GenericRate.lean}.

The type \nolinkurl{Plane} is \nolinkurl{EuclideanSpace} over $\R$ with two
coordinates. The definition \nolinkurl{disk} is the closed unit ball.
\nolinkurl{uniformDisk} is Mathlib's normalized restriction
\nolinkurl{ProbabilityTheory.cond} applied to \texttt{volume disk}; its probability-measure
instance follows from the disk area $\pi>0$.
\texttt{samples N} is the finite product of these probability measures.
The lemma \nolinkurl{samples_rectangle} records the exact product law.

\nolinkurl{hullArea} is the actual Lebesgue measure of the convex hull of
the sample range, converted to a real number and divided by $\pi$.
The project proves finiteness of every such hull measure, its
measurability, and integrability of both area functions. In particular,
the use of real-valued integration does not rely on a default value
for a nonintegrable function. The lemma
\nolinkurl{deficit_eq_expected_volume_deficit} identifies
\texttt{deficit N} with the source's exact normalized expected-volume
deficit.

\nolinkurl{deficit_lower_from_radius} proves \eqref{eq:radius};
\nolinkurl{deficit_lower} specializes the radius and proves the bound.
The generic sequence lemmas have their lower-bound hypothesis
discharged for this exact geometric deficit. The final declarations
establish quadratic divergence, absence of every finite normalized
limit, divergence of the ratio to every positive candidate, and
\nolinkurl{conjecture_00000007788_false}, which negates existence of a
positive constant with the asserted asymptotic equivalence. The
real-power exponent $-2/(2-1)=-2$ is checked explicitly.

\section{Reproduction and verification scope}
The submission includes the complete project, \nolinkurl{Audit.lean}, a
declaration inventory, source hashes, and the portable
\nolinkurl{verify.py} script. With the pinned dependencies and Lean tools
available, run \texttt{python3 verify.py} from the \nolinkurl{lean}
directory. It checks the source hashes and all nine dependency
revisions, rebuilds the authored project, replays every authored Lean
file with warnings treated as errors, and inspects the types and axiom
dependencies of all named declarations. Raw command results and the
verification summary accompany the submission.

There are no auxiliary numerical experiments or finite searches in
this proof: the quantitative and limiting arguments are symbolic.
Only Lean's standard logical axioms
\nolinkurl{propext}, \nolinkurl{Classical.choice}, and \nolinkurl{Quot.sound}
are permitted by the audit. Local verification and independent review
are documented separately from any later maintainer decision.
\end{document}
